\clearpage\pagestyle{otMain}\setlength{\otBlocEspacement}{2.5mm}\fontsize{11}{15}\selectfont\corrigesChapitre{10}{Calculer avec les décimaux}
\begin{blocLivre}{cor-10-1}\begin{corrige}{1}{Une somme posée}\demarche{Compléter $4{,}7$ en $4{,}70$.}
En alignant les virgules : $4{,}70+2{,}85=\res{7{,}55}$. L’estimation $5+3=8$ confirme l’ordre de grandeur.
\end{corrige}\end{blocLivre}
\begin{blocLivre}{cor-10-2}\begin{corrige}{2}{Une différence posée}\demarche{Écrire 10 avec deux chiffres après la virgule.}
On écrit $10{,}00-3{,}48=\res{6{,}52}$. Contrôle : $6{,}52+3{,}48=10$.
\end{corrige}\end{blocLivre}
\begin{blocLivre}{cor-10-3}\begin{corrige}{3}{De tête, sans retenue}\demarche{Ajouter chaque nombre au bon rang.}
L’entier ne touche que la partie entière : $\res{7{,}34}$. Les dixièmes s’ajoutent aux dixièmes : $\res{5{,}54}$. Et $7{,}89-4=\res{3{,}89}$.
\end{corrige}\end{blocLivre}
\begin{blocLivre}{cor-10-4}\begin{corrige}{4}{Par 10, dans les deux sens}\demarche{Déplacer les chiffres d’un rang, pas « ajouter un zéro ».}
Chaque chiffre monte d’un rang : $\res{37}$ et $\res{8{,}5}$. Pour la division, chaque chiffre descend : $\res{4{,}52}$.
\end{corrige}\end{blocLivre}
\begin{blocLivre}{cor-10-5}\begin{corrige}{5}{Un produit en dixièmes}\demarche{Convertir en dixièmes, multiplier, reconvertir.}
24 dixièmes $\times\,3=72$ dixièmes, soit $\res{7{,}2}$. Contrôle : $2{,}4\times3$ est proche de $2\times3=6$ et dépasse 7 de peu.
\end{corrige}\end{blocLivre}
\begin{blocLivre}{cor-10-6}\begin{corrige}{6}{Compléments utiles}\demarche{Compléter d’abord les centièmes, puis les unités.}
$0{,}35+0{,}65=1$ : le complément est $\res{0{,}65}$. $7{,}2+2{,}8=10$ : le complément est $\res{2{,}8}$.
\end{corrige}\end{blocLivre}
\begin{blocLivre}{cor-10-7}\begin{corrige}{7}{Un produit posé}\demarche{Compter les chiffres après la virgule du facteur décimal.}
$345\times6=2\,070$, donc $3{,}45\times6=\res{20{,}70}$, c’est-à-dire $20{,}7$. L’estimation $3{,}5\times6=21$ valide le résultat.
\end{corrige}\end{blocLivre}
\begin{blocLivre}{cor-10-8}\begin{corrige}{8}{Regrouper pour simplifier}\demarche{Chercher un total rond.}
$2{,}6+0{,}4=3$, puis $3+4{,}85=\res{7{,}85}$. L’ordre des termes d’une somme peut changer.
\end{corrige}\end{blocLivre}
\begin{blocLivre}{cor-10-9}\begin{corrige}{9}{Contrôler une soustraction}\demarche{Toujours vérifier une différence par la somme.}
$6{,}35+2{,}35=8{,}70\ne8$ : le résultat est faux. Le bon calcul : $8{,}00-2{,}35=\res{5{,}65}$ ; contrôle : $5{,}65+2{,}35=8$.
\end{corrige}\end{blocLivre}
\begin{blocLivre}{cor-10-10}\begin{corrige}{10}{Rendre la monnaie}\demarche{Avancer de palier en palier jusqu’au billet.}
$7{,}45\to7{,}50$ : 5 centimes ; $7{,}50\to8$ : 50 centimes ; $8\to10$ : 2 euros. La monnaie est $0{,}05+0{,}50+2=\res{2{,}55}$ euros.
\end{corrige}\end{blocLivre}
\begin{blocLivre}{cor-10-11}\begin{corrige}{11}{Le panier de courses}\demarche{Chercher des paires qui font un compte rond.}
La dépense vaut $1{,}20+2{,}85+1{,}15=\res{5{,}20}$ euros ; regrouper $1{,}20+1{,}15=2{,}35$ aide peu, mieux vaut $2{,}85+1{,}15=4$. La monnaie : $10-5{,}20=\res{4{,}80}$ euros.
\end{corrige}\end{blocLivre}
\begin{blocLivre}{cor-10-12}\begin{corrige}{12}{Les sorties du coureur}\demarche{Calculer la donnée manquante avant le total.}
Mardi : $4{,}5-1{,}8=\res{2{,}7}$ km. Total : $4{,}5+2{,}7=\res{7{,}2}$ km. « De moins » conduit ici à une soustraction puis à une addition.
\end{corrige}\end{blocLivre}
\begin{blocLivre}{cor-10-13}\begin{corrige}{13}{Trois carnets identiques}\demarche{Multiplier avant de soustraire.}
$2{,}45\times3=\res{7{,}35}$ euros. Monnaie : $10-7{,}35=\res{2{,}65}$ euros. Estimation : $2{,}5\times3=7{,}5$, proche du prix exact.
\end{corrige}\end{blocLivre}
\begin{blocLivre}{cor-10-14}\begin{corrige}{14}{Le ruban coupé}\demarche{Calculer d’abord la longueur totale coupée.}
Les morceaux prennent $0{,}9\times2=1{,}8$ m. Il reste $2{,}75-1{,}8=\res{0{,}95}$ m. Contrôle : $1{,}8+0{,}95=2{,}75$ m.
\end{corrige}\end{blocLivre}
\begin{blocLivre}{cor-10-15}\begin{corrige}{15}{L’alignement fautif}\demarche{Aligner les virgules, pas les derniers chiffres.}
Tom a ajouté des centièmes à des dixièmes : les rangs étaient décalés. En alignant les virgules : $5{,}60+2{,}14=\res{7{,}74}$. L’estimation $6+2=8$ aurait signalé l’erreur.
\end{corrige}\end{blocLivre}
